\ifdefined\gkver\else\def\gkver{student}\fi
\documentclass[\gkver]{gaokaozhenti}
\begin{document}

\chapter{2016年江苏理综}
%% paperType: 理综物理部分

\begin{enumerate}
\item
%% number: 1
%% paperName: 2016年江苏理综第 1 题 3 分
%% typeId: 1
%% type: 单选题
%% chapter: 力物体的平衡
%% point: 弹力
%% method: 无
%% score: 3
%% degree: 850
%% duplicateId: 0
%% body:
一轻质弹簧原长为 $8\Ucm$，在 $4\UN$ 的拉力作用下伸长了 $2\Ucm$，弹簧未超出弹性限度。则该弹簧的劲度系数为\xzanswer{D}
\fourchoices[answer=D]
{$40\UNmn$}
{$40\UNm$}
{$200\UNmn$}
{$200\UNm$}
%% answer: D
%% memo:
\memoanswer{
由胡克定律得 $F=k\Delta x$，则 $k=\frac{F}{\Delta x}=\frac{4\UN}{0.02\Um}=200\UNm$，故 D 正确。
}

\item
%% number: 2
%% paperName: 2016年江苏理综第 2 题 3 分
%% typeId: 1
%% type: 单选题
%% chapter: 曲线运动
%% point: 斜抛运动
%% method: 无
%% score: 3
%% degree: 750
%% duplicateId: 0
%% body:
有 A、B 两个小球，B 的质量为 A 的两倍。现将它们以相同速率沿同一方向抛出，不计空气阻力。图中①为 A 的运动轨迹，则 B 的运动轨迹是\xzanswer{A}
\onepicture[w=6.5cm]{figs/02.png}
\fourchoices[answer=A]
{①}
{②}
{③}
{④}
%% answer: A
%% memo:
\memoanswer{
因为 A、B 两球只在重力作用下做抛体运动，水平方向上做初速度相同的匀速运动，竖直方向上做初速度和加速度都相同的匀变速运动，故两球的运动轨迹相同，故 A 正确。
}

\item
%% number: 3
%% paperName: 2016年江苏理综第 3 题 3 分
%% typeId: 1
%% type: 单选题
%% chapter: 电场
%% point: 电场线
%% method: 无
%% score: 3
%% degree: 650
%% duplicateId: 0
%% body:
一金属容器置于绝缘板上，带电小球用绝缘细线悬挂于容器中，容器内的电场线分布如图所示。容器内表面为等势面，A、B 为容器内表面上的两点，下列说法正确的是\xzanswer{C}
\onepicture[w=5.0cm]{figs/03.png}
\fourchoices[answer=C]
{A 点的电场强度比 B 点的大}
{小球表面的电势比容器内表面的低}
{B 点的电场强度方向与该处内表面垂直}
{将检验电荷从 A 点沿不同路径到 B 点，电场力所做的功不同}
%% answer: C
%% memo:
\memoanswer{
A 点的电场线比 B 点稀疏，故 A 点的电场强度小于 B 点，故 A 错误；沿着电场线方向电势降低，故小球表面的电势比容器内表面的电势高，故 B 错误；容器内表面是等势面，所以 B 点的电场强度方向与该点容器表面垂直，故 C 正确；由 $W=qU$ 可知，电场力做功与电荷的运动路径无关，故 D 错误。
}

\item
%% number: 4
%% paperName: 2016年江苏理综第 4 题 3 分
%% typeId: 1
%% type: 单选题
%% chapter: 交流电
%% point: 变压器
%% method: 无
%% score: 3
%% degree: 550
%% duplicateId: 0
%% body:
一自耦变压器如图所示，环形铁芯上只绕有一个线圈，将其接在 a、b 间作为原线圈。通过滑动触头取该线圈的一部分，接在 c、d 间作为副线圈。在 a、b 间输入电压为 $U_{1}$ 的交变电流时，c、d 间的输出电压为 $U_{2}$，在将滑动触头从 M 点顺时针旋转到 N 点的过程中\xzanswer{C}
\onepicture[w=6.5cm]{figs/04.png}
\fourchoices[answer=C]
{$U_{2}>U_{1}$，$U_{2}$ 降低}
{$U_{2}>U_{1}$，$U_{2}$ 升高}
{$U_{2}<U_{1}$，$U_{2}$ 降低}
{$U_{2}<U_{1}$，$U_{2}$ 升高}
%% answer: C
%% memo:
\memoanswer{
由图可知，cd 间的线圈为整个线圈的一部分，故其匝数比原线圈 ab 的匝数少，由变压器的电压与匝数的关系 $\frac{U_{1}}{U_{2}}=\frac{n_{1}}{n_{2}}$，可知 $U_{2}<U_{1}$；当触头由 M 点顺时针旋转到 N 点的过程中，cd 间线圈的匝数减少，所以输出电压 $U_{2}$ 降低，故 C 正确。
}

\item
%% number: 5
%% paperName: 2016年江苏理综第 5 题 3 分
%% typeId: 1
%% type: 单选题
%% chapter: 直线运动
%% point: 运动图象
%% method: 无
%% score: 3
%% degree: 550
%% duplicateId: 0
%% body:
小球从一定高度处由静止下落，与地面碰撞后回到原高度再次下落，重复上述运动，取小球的落地点为原点建立坐标系，竖直向上为正方向，下列速度 $v$ 和位置 $x$ 的关系图象中，能描述该过程的是\xzanswer{A}
\fourchoices[answer=A, ispicture=true, h=2.6cm]
{figs/05a.png}
{figs/05b.png}
{figs/05c.png}
{figs/05d.png}
%% answer: A
%% memo:
\memoanswer{
设小球下落的高度为 $H$，依据题意，小球下落过程的位移为 $x_{1}=H-\frac{v^{2}}{2g}$，速度为负值；弹回上升过程的位移为 $x_{2}=H-\frac{v^{2}}{2g}$，速度为正值，故 A 正确。
}

\item
%% number: 6
%% paperName: 2016年江苏理综第 6 题 4 分
%% typeId: 2
%% type: 多选题
%% chapter: 电磁感应
%% point: 法拉第电磁感应定律
%% method: 无
%% score: 4
%% degree: 450
%% duplicateId: 0
%% body:
电吉他中电拾音器的基本结构如图所示，磁体附近的金属弦被磁化，因此弦振动时，在线圈中产生感应电流，电流经电路放大后传送到音箱发生声音，下列说法正确的有\xzanswer{BCD}
\onepicture[w=4.8cm]{figs/06.png}
\fourchoices[answer=BCD]
{选用铜质弦，电吉他仍能正常工作}
{取走磁体，电吉他将不能正常工作}
{增加线圈匝数可以增大线圈中的感应电动势}
{磁振动过程中，线圈中的电流方向不断变化}
%% answer: BCD
%% memo:
\memoanswer{
铜材料不能被磁化，所以选用铜质弦，电吉他不能正常工作，故 A 错误；取走磁体，则没有磁场，不能发生电磁感应现象，电吉他不能正常工作，故 B 正确；由法拉第电磁感应定律，线圈的匝数越多产生的感应电动势越大，故 C 正确；弦振动过程中，线圈的磁场强弱反复变化，由楞次定律，则感应电流的方向不断变化，故 D 正确。
}

\item
%% number: 7
%% paperName: 2016年江苏理综第 7 题 4 分
%% typeId: 2
%% type: 多选题
%% chapter: 曲线运动
%% point: 人造地球卫星
%% method: 无
%% score: 4
%% degree: 450
%% duplicateId: 0
%% body:
如图所示，两质量相等的卫星 A、B 绕地球做匀速圆周运动，用 $R$、$T$、$E_{k}$、$S$ 分别表示卫星的轨道半径、周期、动能、与地心连线在单位时间内扫过的面积。下列关系式正确的有\xzanswer{AD}
\onepicture[w=4.6cm]{\includesvg{figs/07}}
\fourchoices[answer=AD]
{$T_{A}>T_{B}$}
{$E_{kA}>E_{kB}$}
{$S_{A}=S_{B}$}
{$\frac{R_{A}^{3}}{T_{A}^{2}}=\frac{R_{B}^{3}}{T_{B}^{2}}$}
%% answer: AD
%% memo:
\memoanswer{
由开普勒行星运动定律，故 D 正确；单位时间扫过的面积为 $\frac{\pi R^{2}}{T}$，由 $\frac{R_{A}^{3}}{T_{A}^{2}}=\frac{R_{B}^{3}}{T_{B}^{2}}$ 可知 $S_{A}\neq S_{B}$，故 C 错误；因为 $R_{A}>R_{B}$，所以周期 $T_{A}>T_{B}$，故 A 正确；由万有引力提供向心力有 $\frac{GMm}{R^{2}}=m\frac{v^{2}}{R}$，所以卫星 A 的线速度小于 B 的线速度，故 $E_{kA}<E_{kB}$，故 B 错误。
}

\item
%% number: 8
%% paperName: 2016年江苏理综第 8 题 4 分
%% typeId: 2
%% type: 多选题
%% chapter: 电路
%% point: 闭合电路欧姆定律
%% method: 无
%% score: 4
%% degree: 350
%% duplicateId: 0
%% body:
如图所示的电路中，电源电动势为 $12\UV$，内阻为 $2\UO$，四个电阻的阻值已在图中标出。闭合开关 S，下列说法正确的有\xzanswer{AC}
\onepicture[w=6.2cm]{figs/08.png}
\fourchoices[answer=AC]
{路端电压为 $10\UV$}
{电源的总功率为 $10\UW$}
{a、b 间电压的大小为 $5\UV$}
{a、b 间用导线连接后，电路的总电流为 $1\UA$}
%% answer: AC
%% memo:
\memoanswer{
由电路图可知，外电路的总电阻为 $R=\frac{(5+15)(15+5)}{5+15+15+5}\UO=10\UO$，电路的总电流为 $I=\frac{E}{R+r}=\frac{12}{10+2}\UA=1\UA$，路端电压 $U=IR=10\UV$，故 A 正确；电源的总功率 $P_{E}=EI=12\UW$，故 B 错误；两个支路电阻相等，所以各支路中的电流为总电流的一半为 $0.5\UA$，则 a 点右侧的电阻的电势降落为 $0.5\times15\UV=7.5\UV$，同理 b 点右侧的电阻的电势降落为 $0.5\times5\UV=2.5\UV$，故 a、b 间电压大小为 $7.5\UV-2.5\UV=5\UV$，故 C 正确；当 a、b 间用导线连接后，则外电路的总电阻为 $\frac{5\times15}{5+15}\UO+\frac{5\times15}{5+15}\UO=7.5\UO$，总电流为 $I'=\frac{12}{7.5+2}\UA=1.26\UA$，故 D 错误。
}

\item
%% number: 9
%% paperName: 2016年江苏理综第 9 题 4 分
%% typeId: 2
%% type: 多选题
%% chapter: 运动和力
%% point: 传送带
%% method: 无
%% score: 4
%% degree: 250
%% duplicateId: 0
%% body:
如图所示，一只猫在桌边猛地将桌布从鱼缸下拉出，鱼缸最终没有滑出桌面。若鱼缸、桌布、桌面两两之间的动摩擦因数均相等，则在上述过程中\xzanswer{BD}
\onepicture[w=6.0cm]{figs/09.png}
\fourchoices[answer=BD]
{桌布对鱼缸摩擦力的方向向左}
{鱼缸在桌布上的滑动时间和在桌面上的相等}
{若猫增大拉力，鱼缸受到的摩擦力将增大}
{若猫减小拉力，鱼缸有可能滑出桌面}
%% answer: BD
%% memo:
\memoanswer{
鱼缸相对于桌布向左运动，故桌布对鱼缸的摩擦力向右，故 A 错误；因为鱼缸、桌布、桌面两两间的动摩擦因数均相等，则鱼缸受到桌布和桌面的摩擦力相等，对鱼缸的运动而言，先后受到大小相等方向相反的力从静止加速后再减速到零，所以运动的时间相等，故 B 正确；滑动摩擦力和物体间的正压力有关，与拉力无关，所以增大拉力，摩擦力不会变化，故 C 错误；若减小拉力，则桌布运动的加速度将减小，桌布与鱼缸间的相对速度减小，桌布从鱼缸下拉出的时间将增大，鱼缸相对桌面的位移增大，所以鱼缸有可能滑出桌面，故 D 正确。
}

\item
%% number: 10
%% paperName: 2016年江苏理综第 10 题 8 分
%% typeId: 4
%% type: 实验题
%% chapter: 电路
%% point: 电路实验
%% method: 无
%% score: 8
%% degree: 450
%% duplicateId: 0
%% body:
小明同学通过实验探究某一金属电阻的阻值 $R$ 随温度 $t$ 的变化关系。已知该金属电阻在常温下的阻值约 $10\UO$，$R$ 随 $t$ 的升高而增大。实验电路如图所示，控温箱用以调节金属电阻的温度。

实验时闭合 S，先将开关 K 与 1 端闭合，调节金属电阻的温度，分别记下温度 $t_{1}$，$t_{2}$，$\ldots$ 和电流表的相应示数 $I_{1}$，$I_{2}$，$\ldots$。然后将开关 K 与 2 端闭合，调节电阻箱使电流表的示数再次为 $I_{1}$，$I_{2}$，$\ldots$，分别记下电阻箱相应的示数 $R_{1}$，$R_{2}$，$\ldots$。

\onepicture[w=7.5cm]{figs/10a.png}

\begin{enumerate}
	\item
	有以下两电流表，实验电路中应选用\xzanswer{A}

	\twochoices[answer=A]
	{量程 $0\sim100\UmA$，内阻约 $2\UO$}
	{量程 $0\sim0.6\UA$，内阻可忽略}
	\item
	实验过程中，要将电阻箱的阻值由 $9.9\UO$ 调节至 $10.0\UO$，需旋转图中电阻箱的旋钮“a”、“b”、“c”，正确的操作顺序是\tkanswer[2.4cm]{①②③（或①③②）}。

	\begin{steps}
		\item
		将旋钮 a 由“0”旋转至“1”
		\item
		将旋钮 b 由“9”旋转至“0”
		\item
		将旋钮 c 由“9”旋转至“0”
	\end{steps}
	\item
	实验记录的 $t$ 和 $R$ 的数据见下表：

	\begin{center}
	\begin{tabular}{|c|c|c|c|c|c|}
	\hline
	温度 $t/\UdC$ & 20.0 & 40.0 & 60.0 & 80.0 & 100.0 \\
	\hline
	阻值 $R/\UO$ & 9.6 & 10.4 & 11.1 & 12.1 & 12.8 \\
	\hline
	\end{tabular}
	\end{center}

	请根据表中数据，在答题卡的方格纸上作出 $R-t$ 图象。

	\onepicture[w=5.0cm]{\includesvg{figs/10b}}

	由图线求得 $R$ 随 $t$ 的变化关系为 $R=$\tkanswer[2.4cm]{$0.04t+8.8$}$\UO$。
\end{enumerate}
%% answer: 
\jdanswer{
\begin{enumerate}
	\item
	A
	\item
	①②③（或①③②）
	\item
	见图；$0.04t+8.8$（$0.04t+8.6\sim0.04t+9.0$ 都算对）
\end{enumerate}
}
%% memo:
\memoanswer{
（1）电路中电源电动势为 $1.5\UV$，金属电阻约为 $10\UO$，滑动变阻器最大阻值为 $10\UO$，所以电路中的电流的范围约为 $75\UmA$ 到 $150\UmA$，所以量程 $0.6\UA$ 的电流表的指针偏转较小，误差较大，故应选择 A。

（2）为防止电路中电流过大，损坏电表，故实验中电路的电阻的变化应先变大，再变小。

（3）如图所示，则 $R$ 与 $t$ 呈线性关系，在图中取两点，将坐标值代入 $R=kt+R_{0}$，可得 $k=0.04$，$R_{0}=8.8$，所以 $R$ 与 $t$ 的关系式为 $R=0.04t+8.8$。

\onepicture[w=5.5cm]{\includesvg{figs/10_an}}
}

\item
%% number: 11
%% paperName: 2016年江苏理综第 11 题 10 分
%% typeId: 4
%% type: 实验题
%% chapter: 功和能
%% point: DIS研究机械能守恒定律
%% method: 无
%% score: 10
%% degree: 350
%% duplicateId: 0
%% body:
某同学用如图所示的装置验证机械能守恒定律。一根细线系住钢球，悬挂在铁架台上，钢球静止于 A 点，光电门固定在 A 的正下方。在钢球底部竖直地粘住一片宽度为 $d$ 的遮光条。将钢球拉至不同位置由静止释放，遮光条经过光电门的挡光时间 $t$ 可由计时器测出，取 $v=\frac{d}{t}$ 作为钢球经过 A 点时的速度。记录钢球每次下落的高度 $h$ 和计时器示数 $t$，计算并比较钢球在释放点和 A 点之间的势能变化大小 $\Delta E_{\mathrm{p}}$ 与动能变化大小 $\Delta E_{\mathrm{k}}$，就能验证机械能是否守恒。

\onepicture[w=6.0cm]{figs/11a.png}

\begin{enumerate}
	\item
	用 $\Delta E_{\mathrm{p}}=mgh$ 计算钢球重力势能变化的大小，式中钢球下落高度 $h$ 应测量释放时的钢球球心到\xzanswer{B}之间的竖直距离。

	\threechoices[answer=B]
	{钢球在 A 点时的顶端}
	{钢球在 A 点时的球心}
	{钢球在 A 点时的底端}
	\item
	用 $\Delta E_{\mathrm{k}}=\frac{1}{2}mv^{2}$ 计算钢球动能变化的大小，用刻度尺测量遮光条宽度，示数如图所示，其读数为\tkanswer[1.5cm]{1.50}$\Ucm$。某次测量中，计时器的示数为 $0.0100\Us$，则钢球的速度为 $v=$\tkanswer[1.5cm]{1.50}$\Ums$。

	\onepicture[w=4.5cm]{figs/11b.png}
	\item
	下表为该同学的实验结果：

	\begin{center}
	\begin{tabular}{|c|c|c|c|c|c|}
	\hline
	$\Delta E_{\mathrm{p}}/(\times10^{-2}\UJ)$ & 4.892 & 9.786 & 14.69 & 19.59 & 29.38 \\
	\hline
	$\Delta E_{\mathrm{k}}/(\times10^{-2}\UJ)$ & 5.04 & 10.1 & 15.1 & 20.0 & 29.8 \\
	\hline
	\end{tabular}
	\end{center}

	他发现表中的 $\Delta E_{\mathrm{p}}$ 与 $\Delta E_{\mathrm{k}}$ 之间存在差异，认为这是由于空气阻力造成的。你是否同意他的观点？请说明理由。
	\item
	请你提出一条减小上述差异的改进建议。
\end{enumerate}
%% answer: 
\jdanswer{
\begin{enumerate}
	\item
	B
	\item
	1.50（1.49～1.51 都算对）；1.50（1.49～1.51 都算对）
	\item
	不同意，因为空气阻力会造成 $\Delta E_{\mathrm{k}}$ 小于 $\Delta E_{\mathrm{p}}$，但表中 $\Delta E_{\mathrm{k}}$ 大于 $\Delta E_{\mathrm{p}}$
	\item
	分别测出光电门和球心到悬点的长度 $L$ 和 $l$，计算 $\Delta E_{\mathrm{k}}$ 时，将 $v$ 折算成钢球的速度 $v'=\frac{l}{L}v$
\end{enumerate}
}
%% memo:
\memoanswer{
（1）将小球看作质点，释放时从球心开始测量，所以到 A 点时也应该从球心测量。

（2）毫米刻度尺需要估读，速度 $v=\frac{d}{t}=\frac{1.50\times10^{-2}}{0.0100}\Ums=1.50\Ums$。

（3）从实验结果看动能的变化量大于重力势能的变化量，如果是空气阻力的影响，则动能的变化量应小于重力势能的变化量，实验的差异应该是测量小球的速度时测量了小球下端的遮光条的速度，使得速度偏大，从而动能的变化量大于重力势能的变化量。

（4）尽量减小遮光条的尺寸可以减小误差。
}

\item
%% number: 12
%% paperName: 2016年江苏理综第 12 题 4 分
%% typeId: 6
%% type: 计算题
%% chapter: 光学
%% point: 光电效应
%% method: 无
%% score: 4
%% degree: 650
%% duplicateId: 0
%% body:
【选做题】本题包括 A、B、C 三小题，请选定其中两小题，并在相应的答题区域内作答。若多做，则按 A、B 两小题评分。

\textbf{A．}（选修 3-3）（12 分）

\begin{enumerate}
	\item
	在高原地区烧水需要使用高压锅，水烧开后，锅内水面上方充满饱和汽，停止加热，高压锅在密封状态下缓慢冷却，在冷却过程中，锅内水蒸汽的变化情况为\xzanswer{AC}

	\fourchoices[answer=AC]
	{压强变小}
	{压强不变}
	{一直是饱和汽}
	{变为未饱和汽}
	\item
	如图 \ref{2016江苏理综12a} 所示，在斯特林循环的 $p-V$ 图象中，一定质量理想气体从状态 A 依次经过状态 B、C 和 D 后再回到状态 A，整个过程由两个等温和两个等容过程组成，B$\to$C 的过程中，单位体积中的气体分子数目\tkanswer[1.2cm]{不变}（选填“增大”、“减小”或“不变”）；状态 A 和状态 D 的气体分子热运动速率的统计分布图象如图 \ref{2016江苏理综12b} 所示，则状态 A 对应的是\tkanswer[1cm]{①}（选填“①”或“②”）。

	\twopicture[labelA=2016江苏理综12a, labelB=2016江苏理综12b, wA=6.2cm, wB=5.6cm]{figs/12a.png}{figs/12b.png}
	\item
	如图 \ref{2016江苏理综12a} 所示，在 A$\to$B 和 D$\to$A 的过程中，气体放出的热量分别为 $4\UJ$ 和 $30\UJ$。在 B$\to$C 和 C$\to$D 的过程中，气体吸收的热量分别为 $20\UJ$ 和 $12\UJ$。求气体完成一次循环对外界所做的功。
\end{enumerate}

\textbf{B．}（选修 3-4）（12 分）

\begin{enumerate}
	\item
	一艘太空飞船静止时的长度为 $30\Um$，它以 $0.6c$（$c$ 为光速）的速度沿长度方向飞行经过地球，下列说法正确的是\xzanswer{B}

	\fourchoices[answer=B]
	{飞船上的观测者测得该飞船的长度小于 $30\Um$}
	{地球上的观测者测得该飞船的长度小于 $30\Um$}
	{飞船上的观测者测得地球上发来的光信号速度小于 $c$}
	{地球上的观测者测得飞船上发来的光信号速度小于 $c$}
	\item
	杨氏干涉实验证明光的确是一种波，一束单色光投射在两条相距很近的狭缝上，两狭缝就成了两个光源，它们发出的光波满足干涉的必要条件，则两列光的\tkanswer[1.2cm]{频率}相同。如图 \ref{2016江苏理综12c} 所示，在这两列光波相遇的区域中，实线表示波峰，虚线表示波谷，如果放置光屏，在\tkanswer[0.8cm]{C}（选填“A”、“B”或“C”）点会出现暗条纹。

	\onepicture[num=c, label=2016江苏理综12c, w=3.6cm, align=right]{\includesvg{figs/12c}}
	\item
	在上述杨氏干涉实验中，若单色光的波长 $\lambda=5.89\times10^{-7}\Um$，双缝间的距离 $d=1\Umm$，双缝到屏的距离 $l=2\Um$。求第 1 个亮条纹到第 11 个亮条纹的中心之间的距离。
\end{enumerate}

\textbf{C．}（选修 3-5）（12 分）

\begin{enumerate}
	\item
	贝克勒尔在 120 年前首先发现了天然放射现象，如今原子核的放射性在众多领域中有着广泛应用。下列属于放射性衰变的是\xzanswer{A}

	\fourchoices[answer=A]
	{\ce{^{14}_{6}C -> ^{14}_{7}N + ^{0}_{-1}e}}
	{\ce{^{235}_{92}U + ^{1}_{0}n -> ^{131}_{53}I + ^{103}_{39}Y + 2^{1}_{0}n}}
	{\ce{^{2}_{1}H + ^{3}_{1}H -> ^{4}_{2}He + ^{1}_{0}n}}
	{\ce{^{4}_{2}He + ^{27}_{13}Al -> ^{30}_{15}P + ^{1}_{0}n}}
	\item
	已知光速为 $c$，普朗克常数为 $h$，则频率为 $\nu$ 的光子的动量为\tkanswer[1.5cm]{$\frac{h\nu}{c}$}。用该频率的光垂直照射平面镜，光被镜面全部垂直反射回去，则光子在反射前后动量改变量的大小为\tkanswer[1.5cm]{$\frac{2h\nu}{c}$}。
	\item
	几种金属的逸出功 $W_{0}$ 见下表：

	\begin{center}
	\begin{tabular}{|c|c|c|c|c|c|}
	\hline
	金属 & 钨 & 钙 & 钠 & 钾 & 铷 \\
	\hline
	$W_{0}/(\times10^{-19}\UJ)$ & 7.26 & 5.12 & 3.66 & 3.60 & 3.41 \\
	\hline
	\end{tabular}
	\end{center}

	用一束可见光照射上述金属的表面，请通过计算说明哪些能发生光电效应。已知该可见光的波长的范围为 $4.0\times10^{-7}\Um\sim7.6\times10^{-6}\Um$，普朗克常数 $h=6.63\times10^{-34}\UJcdots$。
\end{enumerate}
%% answer: 
\jdanswer{
\begin{enumerate}
	\item
	A：（1）AC；（2）不变，①；（3）$8\UJ$。
	\item
	B：（1）B；（2）频率，C；（3）$L=1.178\times10^{-2}\Um$。
	\item
	C：（1）A；（2）$\frac{h\nu}{c}$，$\frac{2h\nu}{c}$；（3）钠、钾、铷能发生光电效应。
\end{enumerate}
}
%% memo:
\memoanswer{
\textbf{A．}（1）高压锅在密封的状态，所以锅内的气体始终处于饱和状态，故 C 正确；又因为饱和汽压随温度的升高而升高，所以当高压锅逐渐冷却时，温度降低，气压降低，故 A 正确。

（2）气体从状态 B 到状态 C 是等容变化，体积不变，一定质量的气体的分子总数不变，故单位体积内的分子数目不变；气体从状态 D 到状态 A 是等容变化，则压强和温度成正比，故状态 A 的温度小于状态 D 的温度，温度是分子平均动能大小的标志，状态 A 对应的曲线为①。

（3）完成一次循环气体内能不变 $\Delta U=0$，吸收的热量为 $Q=(20+12-4-20)\UJ=8\UJ$，由热力学第一定律 $\Delta U=Q+W$ 得 $W=-8\UJ$，故气体对外做功为 $8\UJ$。

\textbf{B．}（1）由相对论原理可知，飞船上的观察者相对飞船静止，所以其测量的长度不变，飞船以 $0.6c$ 的速度相对于观察者靠近，故地面上的观察者测量的长度小于 $30\Um$，故 A 错误 B 正确；由光速不变原理可知，故 CD 错误。

（2）干涉的必要条件是两列波的频率相同，暗条纹是波峰和波谷相遇的点，故 C 点出现暗条纹。

（3）相邻亮条纹的中心间距 $\Delta x=\frac{l}{d}\lambda$，由题意知，亮条纹的数目 $n=10$，解得 $L=\frac{nl\lambda}{d}$，代入数据得 $L=1.178\times10^{-2}\Um$。

\textbf{C．}（1）天然放射现象分别放射出 $\alpha$ 粒子、$\beta$ 粒子和 $\gamma$ 射线，选项 A 为 $\beta$ 衰变，B 为裂变，C 为轻核聚变，D 为人工核反应，故 A 正确。

（2）光子的动量为 $p=\frac{h}{\lambda}$，光速 $c=\lambda\nu$，所以光子动量 $p=\frac{h\nu}{c}$，动量的变化量 $\Delta p=p_{2}-p_{1}=\frac{h\nu}{c}-\left(-\frac{h\nu}{c}\right)=\frac{2h\nu}{c}$。

（3）光子的能量 $E=\frac{hc}{\lambda}$，取 $\lambda=4.0\times10^{-7}\Um$，则 $E\approx5.0\times10^{-19}\UJ$，由 $E>W_{0}$ 判断，钠、钾、铷能发生光电效应。
}

\item
%% number: 13
%% paperName: 2016年江苏理综第 13 题 15 分
%% typeId: 6
%% type: 计算题
%% chapter: 电磁感应
%% point: 切割磁感线电动势
%% method: 无
%% score: 15
%% degree: 550
%% duplicateId: 0
%% body:
据报道，一法国摄影师拍到“天宫一号”空间站飞过太阳的瞬间。照片中，“天宫一号”的太阳帆板轮廓清晰可见。如图所示，假设“天宫一号”正以速度 $v=7.7\Ukms$ 绕地球做匀速圆周运动，运动方向与太阳帆板两端 M、N 的连线垂直，M、N 间的距离 $L=20\Um$，地磁场的磁感应强度垂直于 $v$、MN 所在平面的分量 $B=1.0\times10^{-5}\UT$，将太阳帆板视为导体。

\begin{enumerate}
	\item
	求 M、N 间感应电动势的大小 $E$；
	\item
	在太阳帆板上将一只“$1.5\UV$、$0.3\UW$”的小灯泡与 M、N 相连构成闭合电路，不计太阳帆板和导线的电阻。试判断小灯泡能否发光，并说明理由；
	\item
	取地球半径 $R=6.4\times10^{3}\Ukm$，地球表面的重力加速度 $g=9.8\Umsq$，试估算“天宫一号”距离地球表面的高度 $h$（计算结果保留一位有效数字）。
\end{enumerate}
\onepicture[w=6.2cm, align=right]{\includesvg{figs/13}}
%% answer: 
\jdanswer{
\begin{enumerate}
	\item
	$E=1.54\UV$
	\item
	不能，因为穿过闭合回路的磁通量不变，不产生感应电流。
	\item
	$h\approx4\times10^{5}\Um$
\end{enumerate}
}
%% memo:
\memoanswer{
（1）由法拉第电磁感应定律得 $E=BLv$，代入数据得 $E=1.54\UV$。

（2）不能，因为穿过闭合回路的磁通量不变，不产生感应电流。

（3）在地球表面，有 $G\frac{Mm}{R^{2}}=mg$，匀速圆周运动的向心力由万有引力提供，有 $G\frac{Mm}{(R+h)^{2}}=m\frac{v^{2}}{R+h}$，解得 $h=\frac{gR^{2}}{v^{2}}-R$，代入数据得 $h\approx4\times10^{5}\Um$（数量级正确都算对）。
}

\item
%% number: 14
%% paperName: 2016年江苏理综第 14 题 16 分
%% typeId: 6
%% type: 计算题
%% chapter: 功和能
%% point: 机械能守恒定律
%% method: 无
%% score: 16
%% degree: 450
%% duplicateId: 0
%% body:
如图所示，倾角为 $\alpha$ 的斜面 A 被固定在水平面上，细线的一端固定于墙面，另一端跨过斜面顶端的小滑轮与物块 B 相连，B 静止在斜面上。滑轮左侧的细线水平，右侧的细线与斜面平行。A、B 的质量均为 $m$。撤去固定 A 的装置后，A、B 均做直线运动。不计一切摩擦，重力加速度为 $g$。求：

\begin{enumerate}
	\item
	A 固定不动时，A 对 B 支持力的大小 $N$；
	\item
	A 滑动的位移为 $x$ 时，B 的位移大小 $s$；
	\item
	A 滑动的位移为 $x$ 时的速度大小 $v_{x}$。
\end{enumerate}
\onepicture[w=7.0cm, align=right]{figs/14.png}
%% answer: 
\jdanswer{
\begin{enumerate}
	\item
	$N=mg\cos\alpha$
	\item
	$s=\sqrt{2(1-\cos\alpha)}\cdot x$
	\item
	$v_{A}=\sqrt{\frac{2gx\sin\alpha}{3-2\cos\alpha}}$
\end{enumerate}
}
%% memo:
\memoanswer{
（1）由受力分析知，支持力的大小 $N=mg\cos\alpha$。

（2）由几何关系得 $s_{x}=x\cdot(1-\cos\alpha)$，$s_{y}=x\cdot\sin\alpha$，且 $s=\sqrt{s_{x}^{2}+s_{y}^{2}}$，联立解得 $s=\sqrt{2(1-\cos\alpha)}\cdot x$。

\onepicture[w=6.0cm]{figs/14_an.png}

（3）B 的下降高度为 $s_{y}=x\cdot\sin\alpha$，由机械能守恒定律得 $mgs_{y}=\frac{1}{2}mv_{A}^{2}+\frac{1}{2}mv_{B}^{2}$，由速度的定义得 $v_{A}=\frac{\Delta x}{\Delta t}$，$v_{B}=\frac{\Delta s}{\Delta t}$，则 $v_{B}=\sqrt{2(1-\cos\alpha)}\cdot v_{A}$，联立解得 $v_{A}=\sqrt{\frac{2gx\sin\alpha}{3-2\cos\alpha}}$。
}

\item
%% number: 15
%% paperName: 2016年江苏理综第 15 题 16 分
%% typeId: 6
%% type: 计算题
%% chapter: 磁场
%% point: 带电粒子在复合场中的运动
%% method: 无
%% score: 16
%% degree: 250
%% duplicateId: 0
%% body:
回旋加速器的工作原理如图 \ref{2016江苏理综15a} 所示，置于真空中的 D 形金属盒半径为 $R$，两盒间狭缝的间距为 $d$，磁感应强度为 $B$ 的匀强磁场与盒面垂直，被加速粒子的质量为 $m$，电荷量为 $+q$，加在狭缝间的交变电压如图 \ref{2016江苏理综15b} 所示，电压值的大小为 $U_{0}$。周期 $T=\frac{2\pi m}{qB}$。一束该粒子在 $t=0\sim\frac{T}{2}$ 时间内从 A 处均匀地飘入狭缝，其初速度视为零。现考虑粒子在狭缝中的运动时间，假设能够出射的粒子每次经过狭缝均做加速运动，不考虑粒子间的相互作用。求：

\begin{enumerate}
	\item
	出射粒子的动能 $E_{k}$；
	\item
	粒子从飘入狭缝至动能达到 $E_{k}$ 所需的总时间 $t_{\text{总}}$；
	\item
	要使飘入狭缝的粒子中有超过 99\% 能射出，$d$ 应满足的条件。
\end{enumerate}
\twopicture[numstyle=arabic, labelprefix={图}, numA=1, numB=2, labelA=2016江苏理综15a, labelB=2016江苏理综15b, wA=6.5cm, wB=6.5cm, align=right]{figs/15a.png}{figs/15b.png}
%% answer: 
\jdanswer{
\begin{enumerate}
	\item
	$E_{m}=\frac{q^{2}B^{2}R^{2}}{2m}$
	\item
	$t_{0}=\frac{\pi BR^{2}+2BRd}{2U_{0}}-\frac{\pi m}{qB}$
	\item
	$d<\frac{\pi mU_{0}}{100qB^{2}R}$
\end{enumerate}
}
%% memo:
\memoanswer{
（1）粒子运动轨迹半径为 $R$ 时，有 $qvB=m\frac{v^{2}}{R}$，且 $E_{m}=\frac{1}{2}mv^{2}$，联立解得 $E_{m}=\frac{q^{2}B^{2}R^{2}}{2m}$。

（2）粒子被加速 $n$ 次达到动能 $E_{m}$，则 $E_{m}=nqU_{0}$，粒子在狭缝间做匀加速运动，设 $n$ 次经过狭缝的总时间为 $\Delta t$，加速度大小为 $a=\frac{qU_{0}}{md}$，由匀加速直线运动的位移公式知 $nd=\frac{1}{2}a\cdot(\Delta t)^{2}$，又 $t_{0}=(n-1)\cdot\frac{T}{2}+\Delta t$，联立解得 $t_{0}=\frac{\pi BR^{2}+2BRd}{2U_{0}}-\frac{\pi m}{qB}$。

（3）经分析，只有在 $0\sim\left(\frac{T}{2}-\Delta t\right)$ 时间内飘入的粒子才能每次均被加速，则所占的比例为 $\eta=\frac{\frac{T}{2}-\Delta t}{\frac{T}{2}}$，由 $\eta>99\%$，解得 $d<\frac{\pi mU_{0}}{100qB^{2}R}$。
}

\end{enumerate}

\end{document}
